\documentclass[12pt]{beamer}
\usepackage[utf8]{inputenc}
\usecolortheme{default}
\usepackage{graphicx}
% \setbeamercovered{transparent}


\newcommand{\prule}{\vspace{-8pt}\rule{50pt}{0.5pt}}




\title[Modale Prädikatenlogik]{Philosophische Modale Prädikatenlogik}
\subtitle{Eine sehr kurze Einführung}
\author{Conrad Friedrich}
\institute[Uni Köln]{Universität zu Köln}
\date{\today}



\begin{document}

\begin{frame}
\titlepage
\end{frame}

\section{Prädikatenlogik}

\begin{frame}
  \Huge Prädikatenlogik \\
  \pause \large ``Baby Logic" Version
\end{frame}

\begin{frame}{Formalisierung}
  \begin{columns}

    \begin{column}[t]{5cm}
      \begin{enumerate}[<+->]
      \item Alle Menschen sind sterblich. \\
      \item Sokrates ist ein Mensch.\\
      \item Also: Sokrates ist sterblich.
      \end{enumerate}
    \end{column}

    \begin{column}[t]{5cm}
      \begin{enumerate}[<+->]
      \item $P$
      \item $Q$
      \item Also: $R$

      \end{enumerate}

      \pause

      \begin{center}
        \includegraphics[height=2cm]{scepticdog.png}
      \end{center}
    \end{column}

  \end{columns}

\pause

  \vspace{10 mm}
  \begin{itemize}
  \item<4-> Mehr Struktur, als wir mit der Aussagenlogik abbilden
    können
  \end{itemize}
\end{frame}

\begin{frame}{Formalisierung}
  
  \begin{columns}

    \begin{column}[t]{5cm}
      \begin{enumerate}[<+->]
      \item Alle Menschen sind sterblich.
      \item Sokrates ist ein Mensch.
      \item Also: Sokrates ist sterblich.
      \end{enumerate}
    \end{column}

    \begin{column}[t]{5cm}
      \begin{enumerate}[<+->]
      \item $\forall x (Mx \to Sx)$
      \item $Ms$
      \item Also: $Ss$

      \end{enumerate}
      \pause
      \begin{center}
         \includegraphics[height=2cm]{cleverdog.png}

      \end{center}

    \end{column}

  \end{columns}

\end{frame}

\begin{frame}{Formalisierung}
  \begin{itemize}
  \item<1-> Es gibt genau einen Gott.
  \item<2-> gdw. Es gibt ein Ding, das Gott ist, und alle anderen
    Dinge sind, falls sie Gott sind, identisch mit diesem Ding.
  \item<3-> gdw. $\exists x(Gx \land \forall y(Gx \to x=y))$.
  \end{itemize}
\end{frame}

\begin{frame}{Vokabular}
  \begin{itemize}
  \item<1-> Variablen
    \begin{itemize}
    \item $\forall \mathbf{x} (M\mathbf{x} \to S\mathbf{x})$
    \end{itemize}
  \item<2-> Konstanten
    \begin{itemize}
    \item $P\mathbf{c}$
    \end{itemize}
  \item<3-> Prädikatensymbole (n-stellig)
    \begin{itemize}
    \item $\forall x (\mathbf{M}x \to \mathbf{S}x)$, $\mathbf{G}xy$
    \end{itemize}
  \item<4-> Konnektive wie in der Aussagenlogik
    \begin{itemize}
    \item $\neg, \to, \land, \lor, \leftrightarrow$
    \end{itemize}
  \item<5-> Quantorsymbole
    \begin{itemize}
    \item $\forall, \exists$
    \end{itemize}
  \item<6-> (Funktionssymbole, Hilfszeichen...)
  \end{itemize}
\end{frame}


\begin{frame}{Grammatik}
  \begin{enumerate}
  \item<1-> Wenn $t_1,...,t_n$ Variablen oder Konstanten sind und $P$
    ein $n$-stelliges Prädikat ist, dann ist $Pt_1,...,t_n$ eine
    (atomare, wohlgeformte) Formel.
    \begin{itemize}
    \item<2-> $Fx$, $Rab$
    \end{itemize}
  \item<3-> Wenn $A$ und $B$ Formeln sind, dann sind auch $\neg A$,
    $A\to B$, $A \land B$, $A \lor B$, $A \leftrightarrow B$ Formeln.
  \item<4-> Wenn A eine Formel ist und $x$ eine Variable, dann sind
    $\forall xA$ und $\exists x A$ Formeln.
  \end{enumerate}
\end{frame}

\begin{frame}{Grammatik: Beispiele}
  Seien $a,b,c$ Konstanten und $P,Q$ Prädikatensymbole. Was wird
  intuitiv mit diesen Formeln ausgedrückt?

  \begin{itemize}
  \item<2-> $Pa$, $Qab$
  \item<3-> $\neg Pa$
  \item<4-> $Qaa$
  \item<5-> $Qab \leftrightarrow Qba$
  \item<6-> $(Qab \land Qbc) \to Qac$
  \end{itemize}

\end{frame}

\begin{frame}{Grammatik: Beispiele}
  Seien $x, y, z$ Variablen, $a,b,c$ Konstanten und $P,Q$
  Prädikatensymbole. Was wird intuitiv mit diesen Formeln ausgedrückt?

  \begin{itemize}
  \item<2-> $Px$, $Qxy$
  \item<3-> $\exists x Px$
  \item<4-> $\forall x (Px \to Qax)$
  \item<5-> $\exists y Qxy$
  \item<6-> $\forall x \exists y Qxy$
  \end{itemize}

  \begin{itemize}
  \item<7-> Variablen, die in einer Formel im Skopus eines Quantors
    stehen, sind \emph{gebunden}, sonst \emph{frei}.
  \item<8-> Formeln, in denen keine freien Variablen vorkommen, heißen
    \emph{Sätze}.
  \end{itemize}
\end{frame}

\begin{frame}{Semantik}
  \begin{itemize}[<+->]
  \item Wann würden wir $Qab$ intuitiv als wahr bezeichnen?

  \item Wenn die Dinge, für die $a,b$ stehen, tatsächlich die
    Eigenschaft haben, die mit $Q$ bezeichnet wird.
  \item $Px$?

  \item Die Variable $x$ hat die Eigenschaft $P$?
  \item Nonsense. Wir können nur \emph{Sätzen} Wahrheitswerte zuordnen
    (zumindest ohne Weiteres.)
  \end{itemize}
\end{frame}

\begin{frame}{Semantik}
  \begin{itemize}[<+->]
  \item Wann würden wir intuitiv $\exists x Px$ als wahr bezeichnen?
  \item Wenn \emph{irgendein} Ding $P$ erfüllt.
  \item $\forall x Fx$?
  \item Wenn \emph{alle} Dinge $F$ erfüllen.
  \end{itemize}
\end{frame}

\begin{frame}{Semantik: Interpretation}
  Eine Interpretation $I$ besteht aus einem Tupel
  $\langle D,v\rangle$.
  \begin{itemize}[<+->]
  \item $D$ ist der (nicht-leere) Gegenstandsbereich, über die
    quantifiziert wird.
  
  \item $v$ ist eine Funktion, so dass
    \begin{itemize}
    \item Wenn $c$ eine Konstante ist, dann ist $v(c) \in D$.
    \item Wenn $P$ ein $n$-stelliges Prädikatensymbol ist, dann ist
      $v(P) \subseteq D^n$.
    \end{itemize}
  \item (Tafel)
  \end{itemize}

\end{frame}

\begin{frame}{Semantik: Wahrheit in einer Interpretation}
  \begin{itemize}[<+->]
  \item $v(Pc_1,...c_n) = 1$ gdw.
    $\langle v(c_1),...,v(c_n)\rangle \in v(P)$, sonst 0.
    \begin{itemize}
    \item $Fa$ ist wahr gdw.\\ $v(a)$ (das Objekt von $a$) in $v(F)$
      (der Extension von $F$) enthalten ist
    \end{itemize}
  \item $v(\neg A)$, $v(A \land B)$ usw. genau wie in der
    Aussagenlogik.
  \end{itemize}

  \begin{itemize}[<+->]
    
  \item $v(\forall xA) = 1$ gdw. \textbf{jedes} Objekt des
    Gegenstandsbereiches $A$ erfüllt, sonst 0.
  \item $v(\exists xA) = 1$ gdw. \textbf{mindestens ein} Objekt des
    Gegenstandsbereiches $A$ erfüllt, sonst 0.
    \begin{itemize}
    \item Eine Formel $A$ gilt in $I$ gdw. $v(I) = 1$.
    \item (Formal unterbestimmt. Was heißt `erfüllen'?)
    \end{itemize}
    % \item $v(\forall xA) = 1$ gdw. \textbf{für alle}
    %   $\partial \in D$
    %   gilt: $v(A_x(c_\partial)) = 1$, sonst 0.
    % \item $v(\exists xA) = 1$ gdw. \textbf{für mindestens ein}
    %   $\partial \in D$ gilt: $v(A_x(c_\partial)) = 1$, sonst 0.

  \end{itemize}

\end{frame}

% \begin{frame}{Semantik V: Wahrheit in einer Interpretation II}
%   \begin{block}{Notation} Sei $\partial \in D$ und $c_\partial$ eine
%     Konstante, sodass $v(c_\partial) = \partial$.  Dann ist
%     $A_x(c_\partial)$ die Formel, die wir erhalten, wenn wir alle
%     Vorkommnisse von $x$ in $A$ durch $c_\partial$ ersetzen.
%     \\
%     \pause
%     \begin{itemize}
%     \item Beispiel: $A$ sei $Px \land Qy$. $A_x(c_\partial)$ ist
%       dann $Pc_\partial \land Qy$.
%     \end{itemize}
%   \end{block}
%   \pause

% \end{frame}

\begin{frame}{Semantik: Interpretation Beispiel}

  Konstanten: $a, b, c$. Prädikatensymbole: $P, Q$. \\
  Sei $I$ gegeben durch: $D = \{\partial_a, \partial_b, \partial_c\}$,\\
  $v(a) = \partial_a$ usw., \\
  $v(P) = \{\partial_a, \partial_b\}$, $v(Q) = \{\langle \partial_a, \partial_a \rangle\, \langle \partial_c, \partial_b \rangle \}$.\\
  Welche der folgenden Formeln gilt in $I$?

  \pause

  \begin{itemize}[<+->]
  \item $Pa \lor Qac$.
  \item $\exists x (Qxx \land Px)$.
  \item $\forall x (Px \to \exists y Qxy)$.
  \item $\forall x (Px \lor Qxy).$
  \end{itemize}

\end{frame}

\begin{frame}{Tableaux-Regeln}

\begin{block}{Notation}
  $A_x(c)$ ist die Formel, die wir erhalten, wenn wir alle freien
  Vorkommnisse von $x$ in $A$ durch $c$ ersetzen.
\end{block}
\pause
\begin{columns}

  \begin{column}[t]{.25 \linewidth}
    \begin{center}
      $\forall x A$ \\
      \prule\\
      $A_x(a)$
    \end{center}
    $a$ ist eine Konstante, die schon vorkam.
  \end{column}
  \pause
  \begin{column}[t]{.25 \linewidth}
    \begin{center}
      $\exists x A$ \\
      \prule\\
      $A_x(c)$
    \end{center}
    $c$ ist eine neue Konstante.
  \end{column}
  \pause
  \begin{column}[t]{.25 \linewidth}
    \begin{center}
      $\neg \exists x A$\\
      \prule\\
      $\forall x \neg A$\\
    \end{center}
  \end{column}
  \pause
  \begin{column}[t]{.25 \linewidth}
    \begin{center}
      $\neg \forall x A$ \\
      \prule\\
      $\exists x \neg A$
    \end{center}
  \end{column}
\end{columns}

\vspace{1em}
\begin{columns}

  \pause

  \begin{column}[t]{.5 \linewidth}
    \begin{itemize}[<+->]
    \item $Pc \vdash \exists x Px$?
    \end{itemize}
  \end{column}

  \begin{column}[t]{.5 \linewidth}
    \begin{itemize}[<+->]
    \item $\forall x \neg Px \vdash \neg \exists x Px$?
    \item $\exists x \neg Px \vdash \neg \forall x Px$?
    \end{itemize}

  \end{column}
\end{columns}

\end{frame}

\begin{frame}
  \Huge Modale Prädikatenlogik\\
  \large Constant Domain
\end{frame}

\begin{frame}{Vokabular und Grammatik}
  \begin{block}{Vokabular}
    Zum Vokabular werden $\Box$ und $\Diamond$ hinzugefügt.
  \end{block}

  \pause
  \begin{block}{Grammatik}
    Wir erweitern die Grammatik der Prädikatenlogik, so dass: \pause
    \begin{itemize}[<+->]
    \item Wenn $A$ eine Formel ist, dann ist auch $\Box A$ eine
      Formel.
    \item Wenn $A$ eine Formel ist, dann ist auch $\Diamond A$ eine
      Formel.
      \begin{itemize}
      \item $\forall x \Box (Px \land Qx) \to \Box \forall x Px$
      \item
        $\Box \Diamond \exists x Px \to \Box \exists x \Diamond (Px
        \lor Qx)$
      \end{itemize}
    \item Restliche Grammatik wie in der klassischen Prädikatenlogik.
    \end{itemize}
  \end{block}

\end{frame}

\begin{frame}{Semantik: Intuitiv}
  Wir wollen die Logik so erweitern, dass: \pause
  \begin{itemize}[<+->]
  \item Ein Gegenstand in einer Situation eine Eigenschaft haben kann,
    aber an einer anderen Welt nicht.
  \item $v(Pa) = 1$ an $w_0$, aber $v(Pb) = 0$ an $w_1$.
  \end{itemize}
\end{frame}

\begin{frame}{Semantik: Interpretation}
  Die Interpretation $I$ wird erweitert, so dass sie aus einem 4-tupel
  $\langle D, W, R, v \rangle$ besteht, wobei \pause
  \begin{itemize}[<+->]
  \item $D$ wie in der klassischen Prädikatenlogik der nicht-leere
    Gegenstandsbereich ist,
  \item $W$ wie in der modalen Aussagenlogik eine Menge möglicher
    Welten ist,
  \item $R$ wie in der modalen Aussagenlogik eine Relation auf $W$ ist
    (d.h. $R \subseteq W \times W$),
  \item $v$ eine Funktion ist, so dass
    \begin{itemize}
    \item Wenn $c$ eine Konstante ist, dann ist $v(c) \in D$.
    \item Wenn $P$ ein $n$-stelliges Prädikatensymbol ist und
      $w \in W$, dann ist $v_w(P) \subseteq D^n$.
    \end{itemize}
  \end{itemize}
\end{frame}

\begin{frame}{Semantik}
  Eine Formel ist nun wahr oder falsch in einer Interpretation
  \emph{an einer Welt}.  Also
  \begin{itemize}[<+->]
  \item $v_\mathbf{w}(Pc_1,...c_n) = 1$ gdw.
    $\langle v(c_1),...,v(c_n)\rangle \in v_\mathbf{w}(P)$, sonst 0.
  \item Ganz analog mit nicht atomaren Formeln:
    \begin{itemize}
    \item $v_w(\neg A) = 1$ gdw. $v_w(A) = 0$.
    \item $v_w(A \land B) = 1$ gdw. $v_w(A) = 1$ und $v_w(B) = 1$,
      sonst 0.
    \item usw.
    \end{itemize}
  \item $v_w(\Box A) = 1$ gdw. \textbf{für alle} $w' \in W$ mit $wRw'$
    gilt: $v_{w'}(A) = 1$.
  \item $v_w(\Diamond A) = 1$ gdw. \textbf{für mindestens ein}
    $w' \in W$ mit $wRw'$ gilt: $v_{w'}(A) = 1$.
  \end{itemize}
\end{frame}

\begin{frame}{Semantik: Beispiel}
  Sei eine Interpretation $I$ gegeben durch: \pause
  \begin{itemize}[<+->]
  \item $D = \{\partial_a, \partial_b, \partial_c\}$
  \item $W = \{w_0, w_1\}$
  \item
    $R = \{\langle w_0, w_0\rangle ,\langle w_0, w_1 \rangle , \langle
    w_1, w_1 \rangle \}$
  \item $v(a) = \partial_a$ usw., $v_{0}(P) = $
  \end{itemize}
  \pause Welche Formeln gelten an der jeweiligen Welt?
  \begin{columns}
    \begin{column}[t]{.5\linewidth}

      \begin{itemize}[<+->]
      \item $\Box Pa$ an $w_1$
      \item $\exists x \Box Px$, $w_0$
      \item $\Box \exists x Px$, $w_0$
      \end{itemize}

    \end{column}

    \begin{column}[t]{.5\linewidth}
      \begin{itemize}[<+->]
      \item $\forall x \Diamond Px$, $w_0$
      \item $\Diamond \forall x Px$, $w_0$
      \end{itemize}
    \end{column}
  \end{columns}

\end{frame}

\begin{frame}{Tableaux: Regeln}

  Alle Regeln der klassischen Prädikatenlogik, relativ zu einer Welt,
  und die Regeln der modalen Aussagenlogik.

  \pause

  \begin{columns}
    \begin{column}[t]{.25 \linewidth}
      \begin{center}
        $\forall x A, i$ \\
        \prule\\
        $A_x(a), i$
      \end{center}
      ($a$ kam schon vor.)
    \end{column}
    
    \begin{column}[t]{.25 \linewidth}
      \begin{center}
        $\exists x A, i$ \\
        \prule\\
        $A_x(c), i$
      \end{center}
      ($c$ ist eine neue Konstante.)
    \end{column}
    
    \begin{column}[t]{.25 \linewidth}
      \begin{center}
        $\neg \exists x A, i$\\
        \prule\\
        $\forall x \neg A, i$\\
      \end{center}
    \end{column}
    
    \begin{column}[t]{.25 \linewidth}
      \begin{center}
        $\neg \forall x A, i$ \\
        \prule\\
        $\exists x \neg A, i$
      \end{center}
    \end{column}
  \end{columns}

  \pause
  
  \begin{columns}
    \begin{column}[t]{.5 \linewidth}
      \begin{center}
        $\Box A, i$ \\
        $irj$ \\
        \prule \\
        $A, j$
      \end{center}
    \end{column}
    \begin{column}[t]{.5 \linewidth}
      \begin{center}
        $\Diamond A, i$\\
        \prule \\
        $irj$ \\
        $A,j$ \\
        ($j$ ist neu.)
      \end{center}

    \end{column}
  \end{columns}

  
\end{frame}



\begin{frame}{Tableaux: Beispiele}

  $K$ ist eine Logik ohne Beschränkung der Relation $R$ der
  Interpretationen.
  \begin{enumerate}[<+->]
  \item $\vdash_K \forall x \Box Px \to \Box \forall x Px$?
  \item $\vdash_K \Box \forall x Px \to \forall x \Box Px$?
  \item $\vdash_K \exists x \Box Px \to \Box \exists x Px$?
  \item $\vdash_K \Box \exists x Px \to \exists x \Box Px$?
  \end{enumerate}


\end{frame}


\begin{frame}
  \Huge Philosophische Probleme und Fragen
\end{frame}

\begin{frame}{Barcan-Formel}

  \begin{itemize}[<+->]
  \item Constant Domain Prädikatenlogik kommt mit einem heftigen
    Commitment:
    \begin{center}
      Barcan Formula (BF): $\forall x \Box Px \to \Box \forall x Px$
    \end{center}
  \item BF ist Theorem von K (und jeder stärkeren Logik).
  \item Wie lässt sich das interpretieren? Nehme an, das Antezedens
    von BF gilt an $w_0$: Alle Dinge an $w_0$ haben die Eigenschaft
    $P$ an allen (zugänglichen) Welten. Dann gilt nach BF auch das
    Konsequenz: An allen Welten gilt $\forall x Px$. D.h. an jeder
    Welt gilt $Px$ für alle \emph{dort} existierenden Dinge.
  \item Ist das plausibel? Was ist, wenn an einer Welt $w_1$ andere
    Dinge existieren als an Welt $w_0$? Dann dürfte
    $\forall x \Box Px$ an $w_0$ nichts über diese Dinge
    implizieren. \end{itemize}
  
\end{frame}

\begin{frame}{Barcan-Formel}

  \begin{itemize}
  \item Äquivalent zur Barcan Formula: 
    \begin{center}
      BF*: $\Diamond \exists x Px \to \exists x \Diamond Px$ 
    \end{center}
  \end{itemize}

  
\end{frame}

\begin{frame}
  Literatur: 
  \begin{itemize}
  \item Fitting, Melvin und Mendelsohn, Richard L. (1998). First Order Modal Logic. Kluwer Academic Publishers.
  \item Garson, James (Spring 2016 Edition). Modal Logic. The Stanford Encyclopedia of Philosophy.
  \item Priest, Graham (2001). Introduction to Non-Classical Logic. Cambridge University Press.
  \end{itemize}

\end{frame}

\end{document}
